\documentclass[12pt]{article}
\usepackage[utf8]{inputenc}
\usepackage[T1]{fontenc}
\usepackage{hyperref}
\title{Thangmi Language Resources for K-12 Curriculum Development}
\author{Compiled by Claude Sonnet 5 (claude-sonnet-5) for Andr\'as Kornai}
\date{}
\begin{document}
\maketitle

\section{Introduction}
Thangmi (ISO 639-3: thf; Glottocode than1259) is an endangered Sino-Tibetan language of Dolakha and Sindhupalchok districts in Nepal, with about 23,000 speakers. Published estimates vary between roughly 23,000 and 30,000. The search found strong academic descriptive work but no evidence of Thangmi-medium school materials.

\section{Primary Educational Resources}
No Thangmi primer, textbook or official curriculum was found in the search results.

\section{Linguistic and Reference Materials}
Mark Turin's 2006 Leiden doctoral dissertation, \emph{A Grammar of the Thangmi Language with an Ethnolinguistic Introduction to the Speakers and Their Culture}, covers phonology, morphology and verb agreement. It also includes glossed texts and a lexicon, and describes the Dolakha dialect:
\url{https://universiteitleiden.nl/en/research/research-output/humanities/a-grammar-of-the-thangmi-language-with-an-ethnolinguistic-introduction-to-the-speakers-and-their-culture}
The full text is hosted in the Leiden repository:
\url{https://scholarlypublications.universiteitleiden.nl/access/item%3A3076307/view}

\section{Religious and Community Resources}
None found in the search results.

\section{Digital and Audio Resources}
None found in the search results.

\section{Curriculum Development Context}
Mark Turin's article ``Mind Your Langs'' argues for mother-tongue instruction as a way to support endangered languages, and notes Thangmi's local plant and animal vocabulary:
\url{https://markturin.sites.olt.ubc.ca/files/2018/07/Turin_Mind_Your_Langs.pdf}
It does not describe an existing Thangmi curriculum.
\end{document}